\documentclass[a4paper,12pt]{article}
\usepackage[T2A]{fontenc}
\usepackage[utf8]{inputenc}
\usepackage[russian]{babel}
\usepackage{amsmath,amssymb,mathtools}
\usepackage{icomma}
\usepackage[margin=19.5mm,top=21mm,bottom=21mm]{geometry}
\usepackage{microtype}
\usepackage{tikz}
\usetikzlibrary{arrows.meta,positioning,shapes.geometric,calc}
\usepackage{enumitem}
\usepackage{tabularx,booktabs,array}
\usepackage{fancyhdr}
\usepackage{parskip}
\usepackage{setspace}
\usepackage{xcolor}
\usepackage{hyperref}
\definecolor{accent}{HTML}{008080}
\definecolor{darkaccent}{HTML}{005558}
\definecolor{textgray}{HTML}{31393A}
\definecolor{muted}{HTML}{647778}
\definecolor{warn}{HTML}{9A4D39}
\hypersetup{colorlinks=true,linkcolor=darkaccent,urlcolor=darkaccent,
  pdfauthor={Лёвин Артём Александрович},
  pdftitle={Чек-лист. Теория вероятностей. 10.11.2026}}
\onehalfspacing
\setlength{\parindent}{0pt}
\setlength{\parskip}{5pt}
\setlist[enumerate]{leftmargin=1.55em,topsep=4pt,itemsep=4pt}
\setlist[itemize]{leftmargin=1.4em,topsep=3pt,itemsep=3pt}
\newcommand{\topic}[1]{\section{#1}\vspace{-2mm}}
\newcommand{\subtopic}[1]{\subsection*{\textcolor{darkaccent}{#1}}\vspace{-1mm}}
\newcommand{\idea}[1]{\textcolor{darkaccent}{\textbf{Главная идея.}}\enspace #1}
\newcommand{\warning}[1]{\textcolor{warn}{\textbf{Обратите внимание.}}\enspace #1}
\newcommand{\answer}[1]{\textcolor{darkaccent}{\textbf{Ответ:}}\enspace #1}
\pagestyle{fancy}
\fancyhf{}
\fancyhead[L]{\footnotesize\color{muted}Теория вероятностей}
\fancyhead[R]{\footnotesize\color{muted}Екатерина Гнедкова}
\fancyfoot[C]{\footnotesize\color{muted}\thepage}
\renewcommand{\headrulewidth}{0.35pt}
\renewcommand{\headrule}{\hbox to\headwidth{\color{accent}\leaders\hrule height \headrulewidth\hfill}}
\renewcommand{\footrulewidth}{0pt}
\setlength{\headheight}{23pt}
\begin{document}
\thispagestyle{empty}
\begin{center}
\vspace*{28mm}
{\large\color{darkaccent} МАТЕМАТИКА}\par
\vspace{21mm}
{\LARGE\color{textgray}Чек-лист}\par
\vspace{12mm}
{\LARGE\color{accent}Теория вероятностей\\[3mm] и случайные события}\par
\vspace{16mm}
{\normalsize Подготовка к ЕГЭ по профильной математике}\par
\vspace{35mm}
{\large Екатерина Гнедкова}\par
\vspace{11mm}
{\normalsize 10 ноября 2026 г.}\par
\vfill
{\small Лёвин Артём Александрович\\[2mm]
эксклюзивно для Екатерины Гнедковой}\par
\vspace{11mm}
\begin{tikzpicture}
  \draw[accent!55,line width=.9pt] (-2.5,0) .. controls (-1.35,.55) and (.8,-.48) .. (2.5,.10);
\end{tikzpicture}
\end{center}
\clearpage

\topic{Теория вероятностей: главное}
\idea{Сначала сформулируйте случайное событие словами. Затем определите множество возможных исходов и выберите формулу, соответствующую условию.}

Порядок повторения: \textbf{равновозможные исходы} $\rightarrow$ \textbf{учёт порядка} $\rightarrow$ \textbf{независимость} $\rightarrow$ \textbf{ровно $k$ успехов} $\rightarrow$ \textbf{дополнение и интервалы}.

\subtopic{1. Равновозможные исходы и учёт порядка}
Если элементарные исходы равновозможны, то
\[
\boxed{P(A)=\frac{N_{\text{благ}}}{N_{\text{всего}}}}.
\]
\textbf{Три выступающие группы (задача Т4.6).} Датская группа должна выступить после шведской и норвежской. Всего $3!=6$ порядков; благоприятны те $2!=2$, в которых датская группа последняя. Вероятность равна $2/6=1/3$.

\textbf{Два броска кубика (Т4.7).} Сумму $8$ дают пять \emph{упорядоченных} пар:
\[
(2;6),\ (3;5),\ (4;4),\ (5;3),\ (6;2).
\]
Из $6\cdot6=36$ исходов подходят пять: $P=5/36\approx0{,}14$.

\textbf{Выбор двух вершин правильного пятиугольника (Т4.12).} Всего пар $\binom52=10$, из них соседних пар, образующих сторону, пять. Поэтому $P=1/2$.

\textbf{Два угла квадрата (Т4.11).} Если первый угол уже выбран, второй выбирают среди трёх оставшихся. Только один угол диагонально противоположен первому. Поэтому вероятность $1/3$.

\textbf{Два близнеца и две равные группы (Т5.12).} После размещения Андрея в одну из групп для Сергея остаётся $25$ мест, из которых $12$ в той же группе. Значит,
\[
P(\text{в одной группе})=\boxed{\frac{12}{25}}.
\]
Этот пример решается условным подсчётом мест: после фиксации группы Андрея у Сергея 25 возможных мест, причём 12 из них подходят. Испытания здесь зависимы; формула Бернулли не требуется.

\subtopic{2. Независимость, произведение, дополнение}
Для независимых событий $A$ и $B$ вероятность их одновременного наступления равна произведению вероятностей. Для любого события $A$ дополнение имеет вероятность $1-P(A)$:
\[
P(A\cap B)=P(A)P(B),\qquad P(\overline A)=1-P(A).
\]
\textbf{Две авторучки (Т8.1).} Каждая неисправна с вероятностью $0{,}1$. При независимости вероятность того, что \emph{обе исправны}, равна
\[
P(\text{обе исправны})=0{,}9\cdot0{,}9=\boxed{0{,}81}.
\]
А вероятность события «\emph{хотя бы одна исправна}» составляет
\[
1-P(\text{обе неисправны})=1-0{,}1^2=0{,}99.
\]
\warning{У этих двух событий разные условия. Всегда определяйте событие словами до записи формулы.}

\textbf{Три последних цифры телефонного номера (Т4.15).} Каждая цифра равновероятно может быть одной из десяти. Нечётных цифр пять, поэтому вероятность, что \emph{все три} нечётны:
\[
P=\left(\frac5{10}\right)^3=\frac18.
\]

\subtopic{3. Ровно $k$ успехов: формула Бернулли}
Для $n$ независимых испытаний с одинаковой вероятностью успеха $p$:
\[
\boxed{P(X=k)=\binom nk p^k(1-p)^{n-k}},\qquad
\binom nk=\frac{n!}{k!(n-k)!}.
\]
\textbf{Пример с тремя независимыми испытаниями.} Для ровно двух успехов при $p=0{,}5$:
\[
P(X=2)=\binom32(0{,}5)^2(0{,}5)=3\cdot0{,}125=0{,}375.
\]
Коэффициент $\binom32=3$ учитывает три расположения двух успехов.

\subtopic{4. Интервалы вероятностей и вложенные события}
Если событие $B$ содержится в событии $A$, то
\[
\boxed{P(A\setminus B)=P(A)-P(B)}.
\]
\textbf{Срок службы фена (Т8.7).} Дано $P(T>1)=0{,}37$ и $P(T>2)=0{,}13$. Требуется $P(1<T\le2)$:
\[
\boxed{P(1<T\le2)=0{,}37-0{,}13=0{,}24}.
\]
\textbf{Масса буханки.} Дано $P(M\le810)=0{,}96$ и $P(M>790)=0{,}93$. Тогда $P(M\le790)=0{,}07$ и
\[
\boxed{P(790<M\le810)=0{,}96-0{,}07=0{,}89}.
\]
\textbf{Диаметр подшипника (Т8.3).} Если $P(|D-67|<0{,}01)=0{,}965$, то для дополнения получаем
\[
\boxed{P(|D-67|\ge0{,}01)=1-0{,}965=0{,}035}.
\]
В формулировке с \emph{строгими} неравенствами $D<66{,}99$ или $D>67{,}01$ из этих данных нельзя однозначно получить ответ: неизвестны вероятности точных граничных значений. Если распределение непрерывное, обе граничные вероятности равны нулю, и ответ также $0{,}035$.

\warning{Сохраняйте указанные в условии строгие и нестрогие границы интервалов. Дополнение к строгому неравенству включает равенство.}

\topic{Итоговая памятка: частые ошибки}
\begin{tabularx}{\linewidth}{@{}p{.33\linewidth}X@{}}
\toprule
\textbf{Ошибка} & \textbf{Как проверить себя}\\
\midrule
Порядок исходов потерян & Два последовательных броска: $(2;6)$ и $(6;2)$ --- разные исходы. При выборе пары вершин порядок не учитывается.\\
\addlinespace[3pt]
«Обе» перепутали с «хотя бы одна» & Для «обе исправны» нужно совместное наступление событий. Для «хотя бы одна исправна» удобно исключить случай «обе неисправны».\\
\addlinespace[3pt]
Забыли число расположений успехов & Для «ровно $k$» при $n$ независимых одинаковых испытаниях используйте коэффициент $\binom nk$.\\
\addlinespace[3pt]
Перепутали «ровно» и «не менее» & «Ровно $k$» --- один член Бернулли; «не менее $k$» --- сумма вероятностей нескольких значений (или дополнение).\\
\addlinespace[3pt]
Неверно вычли вероятности & Для разности $P(A)-P(B)$ проверьте включение событий $B\subseteq A$.\\
\addlinespace[3pt]
Потеряли границы интервала & Внимательно различайте $<$ и $\le$. Дополнение к $X<a$ --- событие $X\ge a$.\\
\bottomrule
\end{tabularx}

\textcolor{darkaccent}{\textbf{Самопроверка:}} любое значение вероятности лежит в отрезке $[0;1]$; смысл итогового события должен полностью соответствовать словам в условии.
\clearpage
\begin{center}
{\Large\color{darkaccent}\textbf{Алгоритм: как выбрать способ подсчёта вероятности}}\par
\vspace{9mm}
\end{center}
\begin{center}
\begin{tikzpicture}[
  >={Stealth[length=3mm]},line width=.8pt,
  b/.style={draw=darkaccent,rounded corners=2pt,text width=10.4cm,align=center,minimum height=10mm,inner sep=5pt,font=\small},
  d/.style={draw=darkaccent,diamond,aspect=3,text width=6cm,align=center,inner sep=4pt,font=\small},
  a/.style={->,accent},
  mini/.style={draw=accent,rounded corners=2pt,align=center,text width=4.75cm,minimum height=13mm,font=\small,inner sep=4pt}]
\node[b] (start) {\textbf{Сначала:} чётко записать событие и все ограничения};
\node[d,below=9mm of start] (type) {Какая структура задачи?};
\node[mini,below left=15mm and -7mm of type] (left) {\textbf{Равновозможные исходы:}\\ $P=N_{\text{благ}}/N_{\text{всего}}$};
\node[mini,below right=15mm and -7mm of type] (right) {\textbf{Независимые испытания:}\\ произведение; при «ровно $k$» --- Бернулли};
\node[b,below=56mm of type] (range) {\textbf{«Хотя бы», «ни разу», диапазон:} проверить дополнение или разность вложенных событий};
\node[b,below=9mm of range] (end) {\textbf{Проверка:} событие сформулировано верно, исходы учтены, $0\le P\le1$};
\draw[a] (start)--(type);
\coordinate (split) at ($(type.south)+(0,-4mm)$);
\draw[accent] (type.south)--(split);
\draw[a] (split) -| (left.north);
\draw[a] (split) -| (right.north);
\coordinate (merge) at ($(range.north)+(0,9mm)$);
\draw[accent] (left.south) |- (merge);
\draw[accent] (right.south) |- (merge);
\draw[a] (merge)--(range.north);
\draw[a] (range)--(end);
\end{tikzpicture}
\end{center}
\clearpage

\clearpage
\begin{center}{\Large\color{darkaccent}\textbf{Тренировочный блок}}\par\end{center}
\vspace{3mm}
\textcolor{muted}{Путь А. 10 заданий по теории вероятностей, от простых к более содержательным. Сначала решите задачи самостоятельно; ответы находятся после блока.}
\begin{enumerate}[label=\textbf{\arabic*.},itemsep=11pt]
\item Правильный кубик бросают один раз. Найдите вероятность выпадения числа, кратного трём.
\item Правильный кубик бросают дважды. Найдите вероятность того, что сумма выпавших очков равна $9$.
\item Три музыкальные группы (датская, шведская и норвежская) выступают в случайном порядке. Найдите вероятность того, что датская группа выступит позже обеих остальных.
\item Из восьми вершин правильного восьмиугольника случайно выбирают две. Найдите вероятность того, что выбранные вершины соединены стороной многоугольника.
\item В классе $28$ человек, в том числе двое близнецов. Учащихся случайно делят на две группы по $14$ человек. Найдите вероятность того, что близнецы окажутся в одной группе.
\item Последние четыре цифры телефонного номера случайны и независимы; каждая цифра от $0$ до $9$ равновероятна. Найдите вероятность того, что все четыре цифры чётные.
\item Две авторучки независимо друг от друга оказываются неисправными каждая с вероятностью $0{,}15$. Найдите вероятность того, что \emph{обе} авторучки исправны.
\item Три устройства выходят из строя независимо друг от друга, каждое с вероятностью $0{,}1$. Найдите вероятность того, что \emph{хотя бы одно} устройство неисправно.
\item В четырёх независимых испытаниях вероятность успеха в каждом равна $0{,}6$. Найдите вероятность того, что успех произойдёт \emph{ровно два} раза.
\item Для массы буханки $M$ известно: $P(M\le820)=0{,}94$ и $P(M>800)=0{,}87$. Найдите вероятность события $800<M\le820$.
\end{enumerate}
\clearpage
\begin{center}{\Large\color{darkaccent}\textbf{Ответы к тренировочному блоку}}\par\end{center}
\vspace{8mm}
\renewcommand{\arraystretch}{1.7}
\begin{tabularx}{\linewidth}{@{}p{1.5cm}X@{}}
\toprule
\textbf{№} & \textbf{Ответ}\\
\midrule
1 & $\dfrac13$.\\
2 & $\dfrac19$.\\
3 & $\dfrac13$.\\
4 & $\dfrac27$.\\
5 & $\dfrac{13}{27}$.\\
6 & $\dfrac1{16}$.\\
7 & $0{,}7225$.\\
8 & $0{,}271$.\\
9 & $0{,}3456$.\\
10 & $0{,}81$.\\
\bottomrule
\end{tabularx}
\vspace{12mm}
\textcolor{darkaccent}{\textbf{Как проверить решение}}\par
Для каждого задания укажите словами искомое событие и объясните выбор формулы. Если исходы равновероятны, проверьте числитель и знаменатель дроби. Если используются независимость или формула Бернулли, проверьте число испытаний, вероятность успеха и необходимое количество успехов. Для интервальной вероятности сначала выпишите вложенные события.
\end{document}
